\subsection{Describing Polynomials (Part 2)} 
\begin{definition}
We can name polynomials by their degree. (These words are \emph{adjectives}.)
\begin{center}
\begin{tabularx}{0.75\linewidth}{c@{\qquad}XX}
\textbf{Degree}&\textbf{Name}&\textbf{Example}\\\hline
0&Constant&$7$\\
1&Linear&$5x-3$\\
2&Quadratic&$x^2-3x+2$\\
3&Cubic&$x^3-5$
\end{tabularx}
\end{center}
\noindent For higher degrees, we just say ``of degree $n$.''
\end{definition}
\begin{definition}
We can name polynomials by their number of terms. (These words are \emph{nouns}.)
\begin{center}
\begin{tabularx}{0.75\linewidth}{c@{\qquad}XXX}
\textbf{Terms}&\textbf{Name}&\multicolumn{2}{l}{\textbf{Examples}}\\\hline
1&Monomial&$x^3$&$2x$\\
2&Binomial&$x-3$&$x^3-8$\\
3&Trinomial&$x^2-3x+2$&$x^5-x^3+x$
\end{tabularx}
\end{center}
\noindent For more terms, we just say ``A polynomial with $n$ terms.''
\end{definition}
\begin{example}Describe each polynomial.
\begin{enumerate}[a.]
\item \makeblankfill{$x^3-7$}
\item \makeblankfill{$x^4+x^5-3$}
\item \makeblankfill{$y^2$}
\item \makeblankfill{$2x^2-5xy+2y^2$}
\item \makeblankfill{$x+y+z+w$}
\item \makeblankfill{$2x-5$}
\item \makeblankfill{$x^3-5x^2+2x-10$}
\item \makeblankfill{$x^4-4x^3y+6x^2y^2-4xy^3+y^4$}
\end{enumerate}
\end{example}